\documentclass[11pt,letterpaper]{article}
\usepackage[utf8]{inputenc}
\usepackage[T1]{fontenc}
\usepackage[spanish,es-noshorthands,es-tabla]{babel}
\usepackage{lmodern}
\usepackage[margin=2.2cm]{geometry}
\usepackage{amsmath,amssymb,amsthm,mathtools}
\usepackage{enumitem}
\usepackage{booktabs,array,tabularx}
\usepackage{multicol}
\usepackage{xcolor}
\usepackage[most]{tcolorbox}
\usepackage{tikz}
\usetikzlibrary{shapes.geometric,arrows.meta,positioning,calc}
\usepackage{fancyhdr}
\usepackage{titlesec}
\usepackage{listings}
\usepackage[hidelinks]{hyperref}

\definecolor{azul}{HTML}{1F4E79}
\definecolor{azulclaro}{HTML}{E8F0F8}
\definecolor{verde}{HTML}{2E7D32}
\definecolor{verdeclaro}{HTML}{EAF5EA}
\definecolor{naranja}{HTML}{B85C00}
\definecolor{naranjaclaro}{HTML}{FFF3E0}
\definecolor{gris}{HTML}{555555}
\definecolor{rojo}{HTML}{B71C1C}

\titleformat{\section}{\Large\bfseries\color{azul}}{\thesection}{0.6em}{}
\titleformat{\subsection}{\large\bfseries\color{azul}}{\thesubsection}{0.5em}{}
\titleformat{\subsubsection}{\normalsize\bfseries\color{gris}}{\thesubsubsection}{0.5em}{}

\setlength{\parindent}{0pt}
\setlength{\parskip}{4pt}
\setlist{itemsep=2pt,topsep=3pt}

\pagestyle{fancy}
\fancyhf{}
\renewcommand{\headrulewidth}{0.4pt}

% ---------- Entornos ----------
\newcounter{ej}
\newtcolorbox{ejercicio}[2][]{enhanced,breakable,colback=white,colframe=azul,
  boxrule=0.6pt,arc=2pt,left=6pt,right=6pt,top=4pt,bottom=4pt,
  title={\textbf{#2}},fonttitle=\small\bfseries,coltitle=white,colbacktitle=azul,#1}
\newtcolorbox{solucion}[2][]{enhanced,breakable,colback=verdeclaro,colframe=verde,
  boxrule=0.6pt,arc=2pt,left=6pt,right=6pt,top=4pt,bottom=4pt,
  title={\textbf{#2}},fonttitle=\small\bfseries,coltitle=white,colbacktitle=verde,#1}
\newtcolorbox{nota}[1][]{enhanced,breakable,colback=naranjaclaro,colframe=naranja,
  boxrule=0.5pt,arc=2pt,left=6pt,right=6pt,top=3pt,bottom=3pt,#1}
\newtcolorbox{formulario}[1][]{enhanced,breakable,colback=azulclaro,colframe=azul,
  boxrule=0.5pt,arc=2pt,left=6pt,right=6pt,top=3pt,bottom=3pt,#1}

\newcommand{\dif}[1]{%
  \ifcase#1\or\textcolor[HTML]{66BB6A}{$\bigstar$}\or\textcolor[HTML]{FFA726}{$\bigstar\bigstar$}\or\textcolor[HTML]{EF5350}{$\bigstar\bigstar\bigstar$}\fi}
\newcommand{\fuente}[1]{\hfill{\footnotesize\textcolor{white!80!azul}{[#1]}}}
\newcommand{\ptst}[1]{{\footnotesize\textcolor{white!80!azul}{(#1 pts)}}}
\newcommand{\pts}[1]{\textcolor{gris}{(#1 pts)}}

% ---------- Probabilidad y RL ----------
\newcommand{\indep}{\perp\!\!\!\perp}
\newcommand{\nindep}{\not\!\perp\!\!\!\perp}
\newcommand{\argmax}{\operatorname*{arg\,max}}
\newcommand{\argmin}{\operatorname*{arg\,min}}
\newcommand{\E}{\mathbb{E}}
\newcommand{\R}{\mathbb{R}}
\tikzset{
  bn/.style={draw,circle,minimum size=0.8cm,inner sep=0pt,font=\small\bfseries,fill=azulclaro},
  bno/.style={bn,fill=naranjaclaro},
  barc/.style={-{Stealth[length=2.2mm]},thick},
  st/.style={draw,rounded corners,minimum width=1.0cm,minimum height=0.8cm,font=\small\bfseries,fill=azulclaro},
  term/.style={st,fill=verdeclaro},
}
% ---------- Lógica ----------
\newcommand{\imp}{\Rightarrow}
\newcommand{\sii}{\Leftrightarrow}
\newcommand{\ent}{\models}
\newcommand{\der}{\vdash}
\newcommand{\p}[1]{\ensuremath{\textit{#1}}}

% ---------- Árboles de juego ----------
\tikzset{
  maxn/.style={draw,regular polygon,regular polygon sides=3,minimum size=0.75cm,inner sep=0pt,fill=azulclaro},
  minn/.style={draw,regular polygon,regular polygon sides=3,shape border rotate=180,minimum size=0.75cm,inner sep=0pt,fill=naranjaclaro},
  chn/.style={draw,circle,minimum size=0.6cm,inner sep=0pt,fill=verdeclaro},
  hoja/.style={draw,rectangle,minimum size=0.5cm,inner sep=1pt,font=\small},
  gnode/.style={draw,circle,minimum size=0.75cm,inner sep=0pt,font=\small\bfseries},
  garc/.style={-{Stealth[length=2mm]},thick},
}

\lstset{basicstyle=\ttfamily\small,columns=fullflexible,frame=single,framerule=0.3pt,rulecolor=\color{gris},
  literate={á}{{\'a}}1 {é}{{\'e}}1 {í}{{\'i}}1 {ó}{{\'o}}1 {ú}{{\'u}}1 {ñ}{{\~n}}1}

% ================= Figuras compartidas (guía y respuestas) — Examen 2 =================

% ---- Red del estudiante ----
\newcommand{\figEst}{%
\begin{tikzpicture}[scale=0.95]
  \node[bn] (D) at (0,2) {D};
  \node[bn] (I) at (2.4,2) {I};
  \node[bn] (C) at (1.2,0.8) {C};
  \node[bn] (E) at (3.6,0.8) {E};
  \node[bn] (R) at (1.2,-0.4) {R};
  \draw[barc] (D) -- (C); \draw[barc] (I) -- (C); \draw[barc] (I) -- (E); \draw[barc] (C) -- (R);
\end{tikzpicture}}

\newcommand{\tabEst}{%
{\small
\begin{tabular}{c}
$P(+d)=0.4$ \qquad $P(+i)=0.3$\\[3pt]
\begin{tabular}{cc|c}
$D$ & $I$ & $P(+c\mid D,I)$\\\hline
$+d$ & $+i$ & 0.7\\ $+d$ & $-i$ & 0.1\\ $-d$ & $+i$ & 0.9\\ $-d$ & $-i$ & 0.4
\end{tabular}\quad
\begin{tabular}{c|c}
$I$ & $P(+e\mid I)$\\\hline $+i$ & 0.8\\ $-i$ & 0.1
\end{tabular}\quad
\begin{tabular}{c|c}
$C$ & $P(+r\mid C)$\\\hline $+c$ & 0.9\\ $-c$ & 0.2
\end{tabular}
\end{tabular}}}

% ---- Grafo de d-separación (8 nodos) ----
\newcommand{\figDsep}{%
\begin{tikzpicture}[scale=0.95]
  \node[bn] (A) at (0,3) {A};
  \node[bn] (B) at (2.4,3) {B};
  \node[bn] (C) at (1.2,1.8) {C};
  \node[bn] (D) at (3.6,1.8) {D};
  \node[bn] (E) at (1.2,0.6) {E};
  \node[bn] (F) at (3.6,0.6) {F};
  \node[bn] (G) at (2.4,-0.6) {G};
  \node[bn] (H) at (2.4,-1.8) {H};
  \draw[barc] (A)--(C); \draw[barc] (B)--(C); \draw[barc] (B)--(D); \draw[barc] (C)--(E);
  \draw[barc] (D)--(F); \draw[barc] (E)--(G); \draw[barc] (F)--(G); \draw[barc] (G)--(H);
\end{tikzpicture}}

% ---- Red de alarma ----
\newcommand{\figAlarma}{%
\begin{tikzpicture}[scale=0.9]
  \node[bn] (B) at (0,1.2) {B}; \node[bn] (E) at (2.4,1.2) {E};
  \node[bn] (A) at (1.2,0) {A};
  \node[bn] (J) at (0,-1.2) {J}; \node[bn] (M) at (2.4,-1.2) {M};
  \draw[barc] (B)--(A); \draw[barc] (E)--(A); \draw[barc] (A)--(J); \draw[barc] (A)--(M);
\end{tikzpicture}}

% ---- Red del aspersor ----
\newcommand{\figAsp}{%
\begin{tikzpicture}[scale=0.9]
  \node[bn] (C) at (1.2,1.2) {C};
  \node[bn] (S) at (0,0) {S}; \node[bn] (R) at (2.4,0) {R};
  \node[bn] (W) at (1.2,-1.2) {W};
  \draw[barc] (C)--(S); \draw[barc] (C)--(R); \draw[barc] (S)--(W); \draw[barc] (R)--(W);
\end{tikzpicture}}

\newcommand{\tabAsp}{%
{\small
\begin{tabular}{c}
$P(+c)=0.5$\qquad $P(+s\mid+c)=0.1$,\ $P(+s\mid-c)=0.5$\qquad $P(+r\mid+c)=0.8$,\ $P(+r\mid-c)=0.2$\\[3pt]
\begin{tabular}{cc|cccc}
$(S,R)$ & & $(+s,+r)$ & $(+s,-r)$ & $(-s,+r)$ & $(-s,-r)$\\\hline
$P(+w\mid S,R)$ & & 0.99 & 0.90 & 0.90 & 0.01
\end{tabular}
\end{tabular}}}

% ---- Red para el orden de eliminación ----
\newcommand{\figOrden}{%
\begin{tikzpicture}[scale=0.9]
  \node[bn] (A) at (0,1.2) {A}; \node[bn] (B) at (2.4,1.2) {B};
  \node[bn] (C) at (1.2,0) {C};
  \node[bn] (D) at (0,-1.2) {D}; \node[bn] (E) at (2.4,-1.2) {E};
  \node[bno] (F) at (1.2,-2.4) {F};
  \draw[barc] (A)--(C); \draw[barc] (B)--(C); \draw[barc] (C)--(D); \draw[barc] (C)--(E);
  \draw[barc] (D)--(F); \draw[barc] (E)--(F);
\end{tikzpicture}}

% ---- Red del servidor ----
\newcommand{\figSrv}{%
\begin{tikzpicture}[scale=0.9]
  \node[bn] (A) at (0,1.2) {A}; \node[bn] (U) at (2.4,1.2) {U};
  \node[bn] (C) at (1.2,0) {C};
  \node[bn] (M) at (0,-1.2) {M}; \node[bn] (Q) at (2.4,-1.2) {Q};
  \draw[barc] (A)--(C); \draw[barc] (U)--(C); \draw[barc] (C)--(M); \draw[barc] (C)--(Q);
\end{tikzpicture}}

\newcommand{\tabSrv}{%
{\small
\begin{tabular}{c}
$P(+a)=0.05$\qquad $P(+u)=0.1$\\[3pt]
\begin{tabular}{cc|c}
$A$ & $U$ & $P(+c\mid A,U)$\\\hline
$+a$ & $+u$ & 0.95\\ $+a$ & $-u$ & 0.80\\ $-a$ & $+u$ & 0.60\\ $-a$ & $-u$ & 0.05
\end{tabular}\quad
\begin{tabular}{c|cc}
$C$ & $P(+m\mid C)$ & $P(+q\mid C)$\\\hline $+c$ & 0.9 & 0.7\\ $-c$ & 0.1 & 0.05
\end{tabular}
\end{tabular}}}

% ---- Red del automóvil (simulacro 1) ----
\newcommand{\figCar}{%
\begin{tikzpicture}[scale=0.9]
  \node[bn] (B) at (0,1.2) {B}; \node[bn] (F) at (2.4,1.2) {F};
  \node[bn] (S) at (1.2,0) {S}; \node[bn] (L) at (-1.2,0) {L};
  \draw[barc] (B)--(S); \draw[barc] (F)--(S); \draw[barc] (B)--(L);
\end{tikzpicture}}

% ---- Red del simulacro 2 ----
\newcommand{\figDdos}{%
\begin{tikzpicture}[scale=0.9]
  \node[bn] (N) at (1.2,1.2) {N};
  \node[bn] (R) at (0,0) {R}; \node[bn] (S) at (2.4,0) {S};
  \node[bn] (W) at (1.2,-1.2) {W}; \node[bn] (T) at (-1.2,-1.2) {T};
  \draw[barc] (N)--(R); \draw[barc] (N)--(S); \draw[barc] (R)--(W); \draw[barc] (S)--(W); \draw[barc] (R)--(T);
\end{tikzpicture}}

% ---- Corredor ----
\newcommand{\figCorredor}{%
\begin{tikzpicture}
  \node[term] (L) at (0,0) {L};
  \node[st] (s1) at (1.6,0) {1};
  \node[st] (s2) at (3.2,0) {2};
  \node[st] (s3) at (4.8,0) {3};
  \node[term] (R) at (6.4,0) {R};
  \node[font=\footnotesize] at (0,-0.75) {entrar: $+2$};
  \node[font=\footnotesize] at (6.4,-0.75) {entrar: $+10$};
  \node[font=\footnotesize] at (3.2,-0.75) {cualquier otro movimiento: $-1$};
\end{tikzpicture}}

% ---- MDP del automóvil (clase) ----
\newcommand{\tabAuto}{%
{\small
\begin{tabular}{cccc}
\toprule
Estado & Acción & Sucesor (prob.) & Recompensa\\\midrule
Cool ($C$) & Slow & $C$ (1) & $+1$\\
Cool ($C$) & Fast & $C$ (0.5), $W$ (0.5) & $+2$\\
Warm ($W$) & Slow & $C$ (0.5), $W$ (0.5) & $+1$\\
Warm ($W$) & Fast & $O$ (1, terminal) & $-10$\\\bottomrule
\end{tabular}}}

\fancyhead[L]{\small\textcolor{gris}{Inteligencia Artificial 2027-I · Guía para el Examen 2}}
\fancyhead[R]{\small\textcolor{gris}{Enunciados}}
\fancyfoot[C]{\small\thepage}

\begin{document}

% ======================= PORTADA =======================
\begin{titlepage}
\centering
\vspace*{1.5cm}
{\color{azul}\rule{\linewidth}{1.2pt}}\\[0.6cm]
{\Huge\bfseries\color{azul} Guía de estudio para el Examen 2}\\[0.3cm]
{\Large Ejercicios a mano de Inteligencia Artificial}\\[0.2cm]
{\large Posgrado en Ciencias e Ingeniería de la Computación, UNAM · 2027-I}\\[0.2cm]
{\large Examen: martes 24 de noviembre de 2026}\\[0.4cm]
{\color{azul}\rule{\linewidth}{1.2pt}}\\[1.2cm]

\begin{minipage}{0.86\linewidth}
\large
\textbf{Cobertura (de agentes probabilísticos en adelante):} probabilidad, independencia y
Bayes; redes bayesianas (factorización, parámetros, d-separación, manto de Markov);
inferencia exacta (enumeración, eliminación de variables, orden); muestreo (previo, rechazo,
verosimilitud con pesos, Gibbs); aprendizaje de parámetros (MLE, Laplace, EM, estructura);
regresión, clasificación logística, softmax y redes; agrupamiento ($k$-NN, $k$-means,
silueta, jerárquico); MDP y aprendizaje por refuerzo (Bellman, evaluación, iteración de valores
y de políticas, Monte Carlo, TD, Q-learning, $\varepsilon$-voraz, aproximación lineal);
aprendizaje profundo (retropropagación, atención, transformers); IA en el contexto social
(métricas de equidad, energía, regulación).

\vspace{0.6cm}
\textbf{Fuentes:} diapositivas 13–25 del curso y apuntes de clase (29 sep – 17 nov);
Russell \& Norvig, \emph{AIMA}, 4.ª ed., caps.\ 12–14, 16–17, 19–23 y 27; Berkeley CS\,188.
Los ejemplos del curso (alarma, aspersor, cadena $R\to T\to L$, automóvil Cool/Warm) se reutilizan
con preguntas nuevas.
\end{minipage}

\vfill
\begin{tabular}{ll}
\dif{1} & Nivel 1 — conceptos y definiciones (respuesta corta)\\[2pt]
\dif{2} & Nivel 2 — procedimientos a mano (tablas, trazas, una o dos iteraciones)\\[2pt]
\dif{3} & Nivel 3 — tipo examen: modelar, ejecutar el algoritmo completo y justificar\\
\end{tabular}

\vspace{1cm}
{\small Las respuestas completas están en el documento complementario
\emph{Guía para el Examen 2 — Respuestas}. Intenta cada ejercicio antes de consultarlas.}
\end{titlepage}

\tableofcontents
\newpage

% ======================= CÓMO USAR =======================
\section*{Cómo usar esta guía}
\addcontentsline{toc}{section}{Cómo usar esta guía}

\begin{itemize}
  \item \textbf{Qué esperar del examen.} En el Examen 1 de este año se pidió desarrollar
  \emph{explícitamente} A* sobre un grafo o minimax a mano; en el de 2025 hubo cuatro preguntas de
  25 puntos que combinaban modelar y ejecutar. El equivalente en este bloque es: \textbf{eliminar
  variables mostrando cada factor, decidir d-separación camino por camino, muestrear con números
  aleatorios dados, contar y normalizar, hacer respaldos de Bellman, actualizar una tabla Q,
  iterar $k$-means y retropropagar en una red pequeña}. Esta guía está construida para eso.
  \item \textbf{Orden.} Parte A (conceptos, \dif{1}), Parte B (procedimientos, \dif{2}), Parte C
  (tipo examen, \dif{3}) y Parte D (dos exámenes simulados de 2 horas, con el formato de
  cuatro preguntas de 25 puntos). Dentro de cada parte se sigue el orden del curso.
  \item \textbf{Convenciones.} $+x$ significa $X$ verdadero y $-x$ falso. Al muestrear con un número
  uniforme $u$, se asigna $+x$ si $u<P(+x\mid\text{padres})$. En $k$-means los empates se asignan al
  grupo de menor índice. En RL, si no se dice otra cosa, $Q_0=0$, los terminales valen 0 y la
  recompensa se cobra \emph{al hacer la transición}. Redondea a 4 cifras, pero arrastra más en
  los cálculos intermedios.
  \item \textbf{Cómo trabajar.} Escribe primero la fórmula que vas a aplicar (por ejemplo,
  $Q(s,a)\leftarrow Q(s,a)+\alpha[r+\gamma\max_{a'}Q(s',a')-Q(s,a)]$) y luego una tabla con un renglón por
  paso. En las tareas y en el Examen 1 se calificó el procedimiento; una tabla ordenada es la
  forma más rápida de mostrarlo.
  \item Un ejercicio de Nivel 3 toma entre 15 y 25 minutos. Cronométrate.
\end{itemize}

\subsection*{Plan de estudio (miércoles 18 → martes 24 de noviembre)}
\begin{center}
\renewcommand{\arraystretch}{1.25}
\begin{tabularx}{\linewidth}{>{\bfseries}l X r}
\toprule
Día & Actividades & Tiempo \\
\midrule
Mié 18 & Formulario (sección 1). Parte A: A1–A4. Parte B: B1–B9 (probabilidad, redes, d-separación, inferencia exacta). & 3–4 h\\
Jue 19 & Parte B: B10–B14 (muestreo y aprendizaje de parámetros). Parte C: C1–C2 (red bayesiana completa a mano). & 3–4 h\\
Vie 20 & Parte A: A5–A6. Parte B: B15–B22 (regresión, logística, softmax, $k$-NN, $k$-means, jerárquico). C3–C5. & 4 h\\
Sáb 21 & Parte A: A7. Parte B: B23–B28 (RL). Parte C: C6–C10 (generalización, $k$-means, MDP, Q-learning, MC vs TD). & 5 h\\
Dom 22 & Parte A: A8–A9. Parte B: B29–B34 (retropropagación, atención, métricas, equidad). C11–C13. & 4 h\\
Lun 23 & Mañana: \textbf{Simulacro D1 cronometrado} (2 h) y corrección. Tarde: \textbf{Simulacro D2} (2 h) y corrección. Noche: repasar los errores y la hoja A4. & 5–6 h\\
Mar 24 & Examen. Antes: releer la hoja A4 y las «trampas frecuentes» del formulario (20 min). & —\\
\bottomrule
\end{tabularx}
\end{center}

\begin{nota}
\textbf{Si tienes poco tiempo} (el proyecto se entrega el 26), prioriza en este orden:
C1 (red bayesiana a mano), B5 (d-separación), B8 (eliminación), B10–B12 (muestreo), C8 (iteración de
valores y políticas), C9 (Q-learning), B20 ($k$-means), B29 (retropropagación) y los dos simulacros.
Son los ejercicios que más se parecen a «desarrollar A* explícitamente» en este bloque.
\end{nota}

\newpage
% ======================= FORMULARIO =======================
\section{Formulario de repaso}

\begin{formulario}
\textbf{Probabilidad.} Marginal $P(x)=\sum_yP(x,y)$. Condicional $P(x\mid y)=P(x,y)/P(y)$.
Producto $P(x,y)=P(x\mid y)P(y)$. Bayes $P(x\mid y)=P(y\mid x)P(x)/P(y)$.
Cadena $P(x_1,\dots,x_n)=\prod_iP(x_i\mid x_{<i})$ (exacta, sin supuestos).
$X\indep Y$ si $P(x,y)=P(x)P(y)$ \emph{para todo} $x,y$; $X\indep Y\mid Z$ si $P(x,y\mid z)=P(x\mid z)P(y\mid z)$ para todo $x,y,z$.
Normalizar: $P(x\mid e)=\alpha\,P(x,e)$, con $\alpha=1/\sum_xP(x,e)$.
\end{formulario}

\begin{formulario}
\textbf{Redes bayesianas.} $P(x_1,\dots,x_n)=\prod_iP(x_i\mid\mathrm{pa}(X_i))$.
Parámetros libres del nodo $X_i$: $(|X_i|-1)\prod_{U\in\mathrm{Pa}(X_i)}|U|$; conjunta completa: $\prod_i|X_i|-1$.
Cada nodo es independiente de sus no descendientes dados sus padres.
\textbf{Manto de Markov} = padres + hijos + otros padres de los hijos; $X$ es independiente del resto dado su manto.

\smallskip
\textbf{d-separación.} $X\indep Y\mid\mathbf Z$ está garantizada si \emph{todo} camino (sin dirección) entre $X$ y $Y$ está bloqueado.
Un camino está activo si \emph{todas} sus tripletas están activas:
\begin{center}\small
\begin{tabular}{lll}
\toprule
Tripleta & Activa cuando & Bloqueada cuando\\\midrule
Cadena $A\to B\to C$ & $B\notin\mathbf Z$ & $B\in\mathbf Z$\\
Causa común $A\leftarrow B\to C$ & $B\notin\mathbf Z$ & $B\in\mathbf Z$\\
Efecto común $A\to B\leftarrow C$ & $B$ \emph{o un descendiente} $\in\mathbf Z$ & ni $B$ ni descendientes en $\mathbf Z$\\\bottomrule
\end{tabular}
\end{center}
Un camino activo \emph{no prueba} dependencia: solo dice que el grafo no garantiza independencia.
\end{formulario}

\begin{formulario}
\textbf{Inferencia exacta.} $P(Q\mid\mathbf e)=\alpha\sum_{\mathbf h}\prod_iP(x_i\mid\mathrm{pa}_i)$.
\emph{Enumeración}: sumar mundos completos (tiempo $O(n\,d^n)$).
\emph{Eliminación de variables}: (1) restringir factores a la evidencia; (2) para cada oculta $H$ en el orden elegido:
unir (multiplicar) todos los factores que mencionan $H$ y sumar $H$; (3) multiplicar lo que queda y normalizar.
El costo lo domina el factor más grande; el orden no cambia la respuesta. Un nodo oculto sin
descendientes en consulta/evidencia es \emph{irrelevante}: $\sum_hP(h\mid\mathrm{pa})=1$.
Inferencia exacta general: NP-difícil (\#P); poliárboles: lineal.
\end{formulario}

\begin{formulario}
\textbf{Muestreo} (orden topológico; $+x$ si $u<P(+x\mid\mathrm{pa})$).
\emph{Previo}: genera de $P(\mathbf X)$; $\hat P(x)=N(x)/N$.
\emph{Rechazo}: descarta muestras incompatibles con $\mathbf e$; acepta $\approx N\,P(\mathbf e)$.
\emph{Verosimilitud con pesos}: fija la evidencia, muestrea el resto, peso $w=\prod_{E_i}P(e_i\mid\mathrm{pa}(E_i))$;
$\hat P(q\mid\mathbf e)=\sum w\,[q]/\sum w$.
\emph{Gibbs}: fija $\mathbf e$, re-muestrea una variable a la vez con
$P(x\mid\text{resto})\propto P(x\mid\mathrm{pa}(X))\prod_{Y\in\mathrm{hijos}}P(y\mid\mathrm{pa}(Y))$ (solo el manto de Markov).
\end{formulario}

\begin{formulario}
\textbf{Aprendizaje de parámetros.} MLE: $\hat P(x\mid u)=N(x,u)/N(u)$ (indefinido si $N(u)=0$).
Laplace: $\hat P(x\mid u)=\dfrac{N(x,u)+\lambda}{N(u)+k\lambda}$, $k=|X|$.
\emph{EM} (oculta $G$): paso E $w_{jg}=P_\theta(g\mid\text{datos}_j)$; paso M: contar con pesos fraccionarios,
$\hat P(g)=\sum_jw_{jg}/m$, $\hat P(x\mid g)=\sum_{j:x_j=x}w_{jg}/\sum_jw_{jg}$. La verosimilitud no decrece; óptimo local.
Estructura: buscar DAG que maximice $\log P(D\mid G,\hat\theta)-\rho\,C(G)$.
\end{formulario}

\begin{formulario}
\textbf{Aprendizaje supervisado.} Regresión: $J(\theta)=\|\mathbf y-X\theta\|^2$, $\nabla J=2X^\mathsf T(X\theta-\mathbf y)$,
ecuaciones normales $X^\mathsf TX\hat\theta=X^\mathsf T\mathbf y$. Para una variable: $\hat\theta_1=\dfrac{\sum(x_i-\bar x)(y_i-\bar y)}{\sum(x_i-\bar x)^2}$, $\hat\theta_0=\bar y-\hat\theta_1\bar x$.
Logística: $\pi=\sigma(\mathbf x^\mathsf T\theta)$, $\sigma(z)=1/(1+e^{-z})$, $\sigma'=\sigma(1-\sigma)$;
$J=-\sum[y\log\pi+(1-y)\log(1-\pi)]$, $\nabla J=X^\mathsf T(\boldsymbol\pi-\mathbf y)$ (convexa). Frontera: $\mathbf x^\mathsf T\theta=0$.
Softmax $p_j=e^{z_j}/\sum_ke^{z_k}$; $\ell=-\log p_{\text{correcta}}$; $\partial\ell/\partial\mathbf z=\mathbf p-\mathbf t$.
Descenso: $\theta\leftarrow\theta-\alpha\nabla J$ (batch, mini-batch, online).
\end{formulario}

\begin{formulario}
\textbf{No supervisado.} $d_p(x,y)=(\sum|x_j-y_j|^p)^{1/p}$; $L_1$, $L_2$, $L_\infty=\max_j|x_j-y_j|$; Hamming = número de coordenadas distintas.
$k$-means: asignar al centroide más cercano, recalcular medias; $J=\sum_j\sum_{x\in C_j}\|x-\mu_j\|^2$ no aumenta; óptimo local.
Silueta $s(i)=\dfrac{b(i)-a(i)}{\max\{a(i),b(i)\}}$ ($a$: media a su grupo; $b$: menor media a otro grupo).
\emph{Linkage}: simple ($\min$), completo ($\max$), promedio, centroide. $k$-NN: votan los $k$ vecinos etiquetados (supervisado).
\end{formulario}

\begin{formulario}
\textbf{MDP y RL.} $G_t=R_{t+1}+\gamma G_{t+1}$.
$V^\pi(s)=\sum_a\pi(a\mid s)\sum_{s'}T(s,a,s')[R(s,a,s')+\gamma V^\pi(s')]$.
$Q^*(s,a)=\sum_{s'}T(s,a,s')[R+\gamma V^*(s')]$, $V^*(s)=\max_aQ^*(s,a)$.
\emph{Iteración de valores}: $V_{k+1}(s)=\max_a\sum_{s'}T[R+\gamma V_k(s')]$ (todos los estados con $V_k$ viejo).
\emph{Iteración de políticas}: evaluar (sistema lineal) $\to$ mejorar $\pi'(s)=\argmax_aQ^\pi(s,a)$ $\to$ repetir hasta que no cambie.
\emph{Monte Carlo}: $V(s)\leftarrow$ promedio de retornos (primera visita / todas las visitas); incremental $V\leftarrow V+\frac1n(G-V)$.
\emph{TD(0)}: $V(s)\leftarrow V(s)+\alpha[r+\gamma V(s')-V(s)]$.
\emph{Q-learning} (fuera de política): $Q(s,a)\leftarrow Q(s,a)+\alpha[r+\gamma\max_{a'}Q(s',a')-Q(s,a)]$; terminal: objetivo $=r$.
\emph{SARSA} (en política): usa $Q(s',a')$ de la acción realmente tomada.
$\varepsilon$-voraz: $\mu(a\mid s)=\varepsilon/m+(1-\varepsilon)\mathbf 1[a\in\mathcal G]/|\mathcal G|$.
Lineal: $\hat Q=\mathbf w^\mathsf T\phi(s,a)$, $\delta=r+\gamma\max_{a'}\hat Q(s',a')-\hat Q(s,a)$, $w_i\leftarrow w_i+\alpha\delta\phi_i$.
\end{formulario}

\begin{formulario}
\textbf{Aprendizaje profundo.} $\mathbf z^{(l)}=W^{(l)}\mathbf a^{(l-1)}+\mathbf b^{(l)}$, $\mathbf a^{(l)}=\phi(\mathbf z^{(l)})$.
Parámetros por capa: $d_l d_{l-1}+d_l$. ReLU$'$ = $1[z>0]$.
Retropropagación: $\boldsymbol\delta^{(L)}=\mathbf p-\mathbf y$ (softmax/sigmoide + entropía cruzada),
$\boldsymbol\delta^{(l)}=(W^{(l+1)\mathsf T}\boldsymbol\delta^{(l+1)})\odot\phi'(\mathbf z^{(l)})$,
$\partial\ell/\partial W^{(l)}=\boldsymbol\delta^{(l)}\mathbf a^{(l-1)\mathsf T}$, $\partial\ell/\partial\mathbf b^{(l)}=\boldsymbol\delta^{(l)}$.
Atención: $A=\mathrm{softmax}_{\text{fila}}(QK^\mathsf T/\sqrt{d_k}+M)$, $O=AV$; máscara causal $M_{ij}=-\infty$ si $j>i$.
Modelo de lenguaje: $p(t_1..t_T)=\prod_ip(t_i\mid t_{<i})$; temperatura $p_j\propto e^{z_j/\tau}$; perplejidad $=e^{\text{pérdida media}}$.
Métricas: exactitud $\frac{TP+TN}{n}$, precisión $\frac{TP}{TP+FP}$, recall/TPR $\frac{TP}{TP+FN}$, especificidad $\frac{TN}{TN+FP}$, FPR $\frac{FP}{FP+TN}$, $F_1=\frac{2TP}{2TP+FP+FN}$.
\end{formulario}

\begin{formulario}
\textbf{Equidad.} Paridad demográfica: $P(\hat Y=1\mid G)$ igual entre grupos. Igualdad de oportunidades:
$P(\hat Y=1\mid Y=1,G)$ (TPR) igual. \emph{Equalized odds}: TPR y FPR iguales. Con tasas base distintas,
estos criterios (y la calibración/PPV) no pueden cumplirse todos a la vez salvo casos triviales.
Energía $E[\mathrm{kWh}]=P[\mathrm W]\cdot t[\mathrm h]/1000$; medir por tarea útil; efecto rebote.
\end{formulario}

\begin{nota}
\textbf{Trampas frecuentes.} (1) Olvidar normalizar o normalizar dos veces. (2) Usar $P(x\mid y)$ por $P(y\mid x)$.
(3) Declarar dependencia porque hay un camino activo. (4) Olvidar que observar un \emph{descendiente} abre un colisionador.
(5) En Laplace, sumar $\lambda$ solo una vez en el denominador (es $k\lambda$). (6) En iteración de valores, usar valores ya
actualizados de la misma iteración. (7) En Q-learning, sumar $\gamma\max Q$ en una transición terminal. (8) Confundir el
$K$ de $k$-means con el $k$ de $k$-NN. (9) En retropropagación, olvidar que ReLU inactiva corta el gradiente. (10) Confundir
precisión (\emph{precision}) con exactitud (\emph{accuracy}).
\end{nota}

\newpage
% ======================= PARTE A =======================
\section{Parte A — Conceptos \dif{1}}

\begin{ejercicio}{A1 · Probabilidad e incertidumbre \dif{1}}
\begin{enumerate}[label=\arabic*.]
\item ¿Por qué un agente lógico puro no basta en un entorno parcialmente observable y estocástico? Da dos razones (pereza, ignorancia).
\item Distingue $P(\mathrm{caries})$, $P(\mathrm{caries}\mid\mathrm{dolor})$ y $P(\mathrm{caries},\mathrm{dolor})$. ¿Cuál suma 1 sobre qué?
\item ¿Cuántas filas y cuántos parámetros libres tiene la conjunta de 8 variables con 3 valores cada una?
\item Verdadero o falso: «si $X\indep Y$ entonces $X\indep Y\mid Z$». Justifica con un ejemplo.
\item ¿Qué afirma la regla de la cadena y qué agrega una red bayesiana sobre ella?
\item Un agente racional bajo incertidumbre elige $\argmax_a\E[U\mid a,\mathbf e]$. ¿Qué dos piezas necesita?
\end{enumerate}
\end{ejercicio}

\begin{ejercicio}{A2 · Redes bayesianas e independencia \dif{1}}
\begin{enumerate}[label=\arabic*.]
\item ¿Qué garantiza la ausencia de una arista entre $X$ y $Y$? ¿Y la presencia?
\item Enuncia la independencia local (de Markov) de un nodo y define el manto de Markov.
\item Explica con un ejemplo cotidiano el fenómeno de «explicar» (\emph{explaining away}).
\item Para cada tripleta (cadena, causa común, efecto común) di qué ocurre al observar el nodo central.
\item ¿Las flechas de una red aprendida de datos indican causalidad? ¿Por qué?
\item ¿Por qué el orden en que se agregan las variables al construir la red cambia su tamaño? (ejemplo: alarma con orden $M,J,A,B,E$).
\end{enumerate}
\end{ejercicio}

\begin{ejercicio}{A3 · Inferencia exacta \dif{1}}
\begin{enumerate}[label=\arabic*.]
\item En $P(Q\mid\mathbf e)$, define consulta, evidencia y variables ocultas. ¿Qué papel tiene $\alpha$?
\item ¿Qué trabajo repetido elimina la eliminación de variables respecto de la enumeración?
\item Define las tres operaciones sobre factores: restringir, unir y eliminar (sumar).
\item ¿Por qué el orden de eliminación cambia el costo pero no el resultado?
\item ¿Qué es un poliárbol y por qué importa? ¿Cuál es la complejidad de la inferencia exacta en general?
\item ¿Qué es una variable irrelevante (\emph{barren}) y cómo se detecta?
\end{enumerate}
\end{ejercicio}

\begin{ejercicio}{A4 · Muestreo y aprendizaje de parámetros \dif{1}}
\begin{enumerate}[label=\arabic*.]
\item ¿Por qué el muestreo por rechazo es ineficiente cuando $P(\mathbf e)$ es pequeña? ¿Cuántas muestras útiles esperas de $N$?
\item En verosimilitud con pesos, ¿por qué no basta contar las muestras sin peso? ¿Qué problema aparece si la evidencia está «abajo» y es rara?
\item ¿Qué necesita Gibbs para re-muestrear una variable? ¿Por qué sus muestras están correlacionadas y qué es el calentamiento?
\item ¿Qué se maximiza en MLE? ¿Por qué con datos completos basta «contar y normalizar» cada fila de cada CPT?
\item ¿Para qué sirve Laplace? ¿Qué ocurre con $\lambda\to\infty$?
\item ¿Qué hacen los pasos E y M de EM? ¿Qué garantiza y qué no?
\item ¿Por qué la verosimilitud sola no sirve para elegir la estructura de una red?
\end{enumerate}
\end{ejercicio}

\begin{ejercicio}{A5 · Aprendizaje supervisado \dif{1}}
\begin{enumerate}[label=\arabic*.]
\item Distingue clasificación de regresión. ¿Qué significa «lineal» en regresión polinomial?
\item ¿Cuándo existe $(X^\mathsf TX)^{-1}$? ¿Qué se hace si no existe?
\item Define sobreajuste y subajuste en términos de error de entrenamiento y de prueba.
\item ¿Por qué la regresión logística minimiza entropía cruzada y no error cuadrático? ¿Qué propiedad tiene su pérdida?
\item ¿Por qué un perceptrón (o la logística) no resuelve XOR? ¿Qué agrega una capa oculta?
\item Compara descenso batch, mini-batch y en línea.
\item ¿Qué afirma (y qué \emph{no}) el teorema de aproximación universal?
\item ¿Qué aporta un proceso gaussiano frente a una predicción puntual?
\end{enumerate}
\end{ejercicio}

\begin{ejercicio}{A6 · Aprendizaje no supervisado \dif{1}}
\begin{enumerate}[label=\arabic*.]
\item ¿Por qué hay que escalar las características antes de usar distancias?
\item ¿Qué minimiza $k$-means y por qué converge? ¿Converge al óptimo global?
\item Interpreta una silueta de $0.8$, $0$ y $-0.3$.
\item ¿Qué tipo de grupos favorece el \emph{linkage} simple y cuál el completo?
\item ¿Cuándo preferirías DBSCAN a $k$-means?
\item ¿$k$-NN es supervisado o no supervisado? ¿Qué es $k$ y qué es $K$?
\end{enumerate}
\end{ejercicio}

\begin{ejercicio}{A7 · MDP y aprendizaje por refuerzo \dif{1}}
\begin{enumerate}[label=\arabic*.]
\item Lista los componentes de un MDP. ¿Qué es la propiedad de Markov?
\item ¿Por qué se usa descuento $\gamma<1$? Acota $|G_t|$ si $|R|\le R_{\max}$.
\item Distingue $V^\pi$, $V^*$, $Q^\pi$, $Q^*$ y política. ¿Por qué con $Q^*$ no hace falta el modelo para actuar?
\item ¿Por qué iteración de valores converge? (contracción con factor $\gamma$).
\item Compara iteración de valores con iteración de políticas.
\item ¿Qué diferencia hay entre planificación (MDP conocido) y aprendizaje por refuerzo? ¿Basado en modelo vs.\ libre de modelo?
\item Compara Monte Carlo y TD: ¿qué necesita cada uno, qué sesgo y qué varianza tienen?
\item ¿Por qué Q-learning es «fuera de política» y SARSA «en política»?
\item ¿Por qué hay que explorar? ¿Qué condiciones requiere la convergencia de Q-learning tabular?
\item ¿Qué problema resuelve la aproximación de funciones y qué se pierde? ¿Qué agrega DQN (repetición de experiencias, red objetivo)?
\end{enumerate}
\end{ejercicio}

\begin{ejercicio}{A8 · Aprendizaje profundo \dif{1}}
\begin{enumerate}[label=\arabic*.]
\item ¿Por qué componer capas lineales sin activación no aumenta la expresividad?
\item ¿Qué problema de la sigmoide reduce ReLU? ¿Qué es una neurona muerta?
\item ¿Qué calcula la retropropagación y por qué es eficiente?
\item ¿Qué es un \emph{embedding}? ¿Qué mide la similitud coseno?
\item En atención, ¿qué papel juegan consultas, claves y valores? ¿Por qué se divide entre $\sqrt{d_k}$? ¿Para qué es la máscara causal?
\item Distingue preentrenamiento, ajuste por instrucciones y RLHF.
\item ¿Qué efecto tiene la temperatura en la decodificación?
\item Explica el compromiso sesgo–varianza.
\end{enumerate}
\end{ejercicio}

\begin{ejercicio}{A9 · IA en el contexto social \dif{1}}
\begin{enumerate}[label=\arabic*.]
\item Nombra cuatro fuentes de discriminación algorítmica (antes, durante y después del modelo).
\item ¿Por qué eliminar el atributo protegido no garantiza equidad?
\item Define paridad demográfica, igualdad de oportunidades y \emph{equalized odds}. ¿Por qué pueden ser incompatibles?
\item ¿Qué es la inyección de instrucciones (\emph{prompt injection}) y qué controles la mitigan en un agente?
\item Describe la pirámide de riesgo del Reglamento de IA de la UE con un ejemplo por nivel.
\item Distingue potencia (W) de energía (kWh). ¿Qué es el efecto rebote?
\item Distingue exposición de una ocupación a la IA y sustitución. ¿Qué es la deuda cognitiva?
\item Enumera las «cuatro D» para colaborar con IA.
\end{enumerate}
\end{ejercicio}

\newpage
% ======================= PARTE B =======================
\section{Parte B — Procedimientos a mano \dif{2}}

\subsection{Probabilidad}

\begin{ejercicio}{B1 · Conjunta, marginales e independencia \dif{2}\fuente{Clase 29 sep}}
La conjunta de $L$ (lluvia), $T$ (tráfico) y $P$ (paraguas) es:
\begin{center}\small
\begin{tabular}{ccc|c@{\qquad}ccc|c}
$L$ & $T$ & $P$ & $P(\cdot)$ & $L$ & $T$ & $P$ & $P(\cdot)$\\\hline
$+$ & $+$ & $+$ & 0.216 & $-$ & $+$ & $+$ & 0.014\\
$+$ & $+$ & $-$ & 0.024 & $-$ & $+$ & $-$ & 0.126\\
$+$ & $-$ & $+$ & 0.054 & $-$ & $-$ & $+$ & 0.056\\
$+$ & $-$ & $-$ & 0.006 & $-$ & $-$ & $-$ & 0.504
\end{tabular}
\end{center}
(a) Calcula $P(+t)$, $P(+p)$, $P(+t,+p)$. ¿Son $T$ y $P$ independientes?
(b) Calcula $P(+l\mid+t)$ y $P(+l\mid+t,+p)$. Interpreta.
(c) Verifica si $T\indep P\mid L$ (ambos valores de $L$).
(d) ¿Qué red de 3 nodos es consistente con lo anterior? ¿Cuántos parámetros ahorra frente a la conjunta?
\end{ejercicio}

\begin{ejercicio}{B2 · Bayes con dos pruebas \dif{2}\fuente{AIMA 12.x, estilo}}
Una enfermedad tiene prevalencia del 2\,\%. Una prueba tiene sensibilidad $0.95$ y tasa de falsos positivos $0.08$.
(a) Calcula $P(\text{enf}\mid+)$.
(b) Se aplica una segunda prueba igual, condicionalmente independiente dada la enfermedad. Calcula
$P(\text{enf}\mid+,+)$ de dos formas: con la conjunta y actualizando secuencialmente.
(c) Calcula $P(\text{enf}\mid+,-)$.
(d) Dibuja la red bayesiana que justifica el supuesto de (b).
\end{ejercicio}

\begin{ejercicio}{B3 · Factorización y conteo de parámetros \dif{2}}
Variables con dominios $|A|=2$, $|B|=3$, $|C|=2$, $|D|=4$, $|E|=2$, $|F|=3$ y aristas $A\to C$, $B\to C$, $C\to D$, $C\to E$, $E\to F$.
(a) Escribe la factorización. (b) Cuenta parámetros libres por nodo y en total. (c) Compara con la conjunta.
(d) Si se agrega la arista $A\to F$, ¿cuántos parámetros más se requieren?
\end{ejercicio}

\subsection{Redes bayesianas e independencia}

\begin{ejercicio}{B4 · Red del estudiante: mundos y marginales \dif{2}}
\begin{center}\figEst\qquad\tabEst\end{center}
$D$: curso difícil; $I$: estudiante inteligente; $C$: calificación alta; $E$: examen de admisión alto; $R$: buena carta de recomendación.
(a) Escribe la factorización y cuenta parámetros (red vs.\ conjunta).
(b) Calcula $P(-d,+i,+c,-e,+r)$.
(c) Calcula $P(+c)$ sumando sobre $D$ e $I$.
(d) Sin calcular, ¿$D\indep I$? ¿$D\indep I\mid R$? ¿$D\indep E$? ¿$R\indep E\mid C$?
\end{ejercicio}

\begin{ejercicio}{B5 · d-separación camino por camino \dif{2}\fuente{Estilo diapositivas 27–29}}
\begin{center}\figDsep\end{center}
Para cada consulta, lista los caminos simples relevantes, marca el estado de cada tripleta y concluye si la
independencia está \emph{garantizada}:
\begin{multicols}{2}
\begin{enumerate}[label=(\alph*)]
\item $A\indep B$
\item $A\indep B\mid H$
\item $A\indep D$
\item $A\indep D\mid G$
\item $A\indep D\mid G,B$
\item $C\indep F\mid B$
\item $E\indep F\mid G,D$
\item $A\indep H\mid E$
\item $A\indep F\mid C$
\item $B\indep G\mid E,F$
\end{enumerate}
\end{multicols}
(k) Da el manto de Markov de $C$, $D$, $F$ y $G$.
\end{ejercicio}

\begin{ejercicio}{B6 · Alarma: explicar a la otra causa \dif{2}\fuente{Clase, red de alarma}}
\begin{center}\figAlarma\end{center}
Usa las CPT del curso: $P(+b)=0.001$, $P(+e)=0.002$, $P(+a\mid B,E)=0.95,0.94,0.29,0.001$ para
$(+b,+e),(+b,-e),(-b,+e),(-b,-e)$.
(a) Calcula $P(+a)$. (b) Calcula $P(+b\mid+a)$ y $P(+e\mid+a)$.
(c) Calcula $P(+e\mid+a,+b)$ y $P(+e\mid+a,-b)$. Explica el resultado en una línea.
(d) ¿Por qué $J$ y $M$ no intervienen en (b) y (c)?
\end{ejercicio}

\subsection{Inferencia exacta}

\begin{ejercicio}{B7 · Enumeración en una cadena \dif{2}\fuente{Clase 13 oct, diap.\ 18–22}}
$R\to T\to L$ con $P(+r)=0.1$; $P(+t\mid+r)=0.8$, $P(+t\mid-r)=0.1$; $P(+l\mid+t)=0.3$, $P(+l\mid-t)=0.1$.
(a) Construye la tabla de los 8 mundos. (b) Calcula $P(R\mid+l)$ y $P(T\mid+l)$ por enumeración.
(c) Calcula $P(L\mid+r)$. ¿Por qué aquí no hace falta normalizar al final?
\end{ejercicio}

\begin{ejercicio}{B8 · Eliminación de variables, factor por factor \dif{2}}
En la red del estudiante (B4) calcula $P(I\mid+r)$ por eliminación de variables con orden $D$, $C$.
(a) Lista los factores iniciales tras restringir a $+r$. (b) Justifica por qué $P(E\mid I)$ puede descartarse.
(c) Construye $f_1(C,I)=\sum_dP(d)P(C\mid d,I)$ (4 entradas). (d) Construye $f_2(I)=\sum_cf_1(c,I)P(+r\mid c)$.
(e) Multiplica por $P(I)$ y normaliza.
\end{ejercicio}

\begin{ejercicio}{B9 · Orden de eliminación y tamaño de factores \dif{2}\fuente{Clase 13 oct, diap.\ 29}}
\begin{center}\figOrden\end{center}
Variables booleanas; se observa $+f$ y se consulta $P(A\mid+f)$. Para cada orden lista, en cada paso, los factores
que se unen, el alcance del producto y el factor nuevo: (i) $B,D,E,C$; (ii) $C,B,D,E$; (iii) $D,E,C,B$.
¿Cuál es el factor más grande de cada orden (en entradas)? ¿Cuál recomiendas y por qué?
\end{ejercicio}

\subsection{Muestreo}

\begin{ejercicio}{B10 · Muestreo previo y rechazo con números dados \dif{2}\fuente{Clase 15 oct}}
\begin{center}\figAsp\qquad\tabAsp\end{center}
Genera 5 muestras en el orden $C,S,R,W$ usando, para cada muestra, los cuatro uniformes:
\begin{center}\small
\begin{tabular}{c|cccc}
Muestra & $u_C$ & $u_S$ & $u_R$ & $u_W$\\\hline
1 & 0.31 & 0.62 & 0.12 & 0.45\\
2 & 0.77 & 0.08 & 0.91 & 0.50\\
3 & 0.18 & 0.35 & 0.66 & 0.97\\
4 & 0.52 & 0.44 & 0.15 & 0.03\\
5 & 0.05 & 0.06 & 0.70 & 0.88
\end{tabular}
\end{center}
(a) Estima $P(+w)$ (exacto: $0.65$). (b) Con rechazo, estima $P(+c\mid+w)$ (exacto: $0.5746$).
(c) ¿Puedes estimar $P(+c\mid-r,-w)$ con estas muestras? ¿Por qué?
\end{ejercicio}

\begin{ejercicio}{B11 · Verosimilitud con pesos \dif{2}}
Misma red. Evidencia: $+s,+w$. Consulta: $P(+c\mid+s,+w)$. Se muestrean solo $C$ y $R$, con los pares
$(u_C,u_R)$: $(0.21,0.40)$, $(0.64,0.10)$, $(0.47,0.93)$, $(0.88,0.75)$, $(0.12,0.85)$.
(a) Para cada muestra da $C$, $R$ y el peso $w=P(+s\mid c)\,P(+w\mid+s,r)$.
(b) Estima $P(+c\mid+s,+w)$. (c) Calcula el valor exacto y compara. (d) ¿Por qué el peso incluye $P(+s\mid c)$?
\end{ejercicio}

\begin{ejercicio}{B12 · Dos barridos de Gibbs \dif{2}}
Misma red. Evidencia $+w$. Estado inicial $(+c,-s,-r)$. Orden de actualización $C,S,R$, dos barridos, con uniformes
$0.60,\ 0.20,\ 0.30,\ 0.90,\ 0.55,\ 0.40$ (en ese orden).
(a) Escribe, para cada variable, la distribución condicional usando solo su manto de Markov.
(b) Ejecuta los dos barridos mostrando pesos, probabilidad normalizada y nuevo valor.
(c) ¿Qué estimación de $P(+c\mid+w)$ dan los dos estados? ¿Por qué no es confiable?
\end{ejercicio}

\subsection{Aprendizaje de parámetros}

\begin{ejercicio}{B13 · MLE y Laplace con dos padres \dif{2}\fuente{Clase 20 oct, estilo}}
Red $\text{Clima}\to\text{Bici}\leftarrow\text{Día}$, con Clima $\in\{S,N,Ll\}$ (sol, nublado, lluvia), Día $\in\{L,F\}$ (laboral, fin de semana), Bici $\in\{\text{sí},\text{no}\}$.
\begin{center}\small
\begin{tabular}{c|cccccccccccc}
\# & 1&2&3&4&5&6&7&8&9&10&11&12\\\hline
Clima & S&S&S&S&N&N&N&Ll&Ll&S&N&Ll\\
Día & L&L&F&L&L&F&L&L&L&F&F&L\\
Bici & sí&sí&sí&no&sí&no&no&no&no&sí&sí&sí
\end{tabular}
\end{center}
(a) Estima $P(\text{Clima})$ por MLE y con Laplace $\lambda=1$.
(b) Estima por MLE cada fila $P(\text{Bici}=\text{sí}\mid\text{Clima},\text{Día})$. ¿Qué fila queda indefinida?
(c) Repite (b) con Laplace $\lambda=1$. (d) ¿Cuántos parámetros libres tiene la red?
\end{ejercicio}

\begin{ejercicio}{B14 · Una iteración de EM \dif{2}\fuente{Clase 20 oct, diap.\ 19–21}}
Red $G\to R_1$, $G\to R_2$ con $G\in\{a,b\}$ oculta y $R_1,R_2\in\{0,1\}$ (le gustó o no). Datos:
$(1,1),(1,0),(0,0),(1,1),(0,1)$. Parámetros iniciales: $P(a)=0.5$; $P(R_1{=}1\mid a)=0.8$, $P(R_1{=}1\mid b)=0.3$;
$P(R_2{=}1\mid a)=0.7$, $P(R_2{=}1\mid b)=0.4$.
(a) Paso E: calcula $w_j=P(a\mid r_{1j},r_{2j})$ para cada registro.
(b) Paso M: reestima los cinco parámetros.
(c) Calcula la log-verosimilitud $\sum_j\log P(r_{1j},r_{2j})$ antes y después. ¿Qué observas?
\end{ejercicio}

\subsection{Aprendizaje supervisado}

\begin{ejercicio}{B15 · Regresión lineal y descenso de gradiente \dif{2}\fuente{Clase 22 oct}}
Datos $(x,y)$: $(0,1)$, $(1,3)$, $(2,4)$, $(3,8)$; modelo $\hat y=\theta_0+\theta_1x$, $J=\sum(y-\hat y)^2$.
(a) Escribe $X$, $X^\mathsf TX$ y $X^\mathsf T\mathbf y$; resuelve las ecuaciones normales.
(b) Calcula predicciones, residuos y $J(\hat\theta)$. Predice para $x=5$.
(c) Desde $\theta=(0,0)$, da dos pasos de descenso con $\alpha=0.01$ y reporta $J$ tras cada paso.
\end{ejercicio}

\begin{ejercicio}{B16 · Regresión logística: pérdida, gradiente y un paso \dif{2}\fuente{Clase 27 oct}}
Tres ejemplos con $\mathbf x=(1,x_1,x_2)$: $(1,2,1)$ con $y=1$; $(1,0,2)$ con $y=0$; $(1,1,0)$ con $y=1$. Parámetros $\theta=(0,1,-1)$.
(a) Calcula $z_i$, $\pi_i$ y la entropía cruzada total. (b) Calcula $\nabla J=X^\mathsf T(\boldsymbol\pi-\mathbf y)$.
(c) Da un paso con $\alpha=0.5$ y recalcula la pérdida. (d) Escribe la frontera de decisión antes y después.
\end{ejercicio}

\begin{ejercicio}{B17 · Softmax, entropía cruzada y temperatura \dif{2}}
Logits $\mathbf z=(2,1,-1)$, clase correcta = 2.
(a) Calcula $\mathbf p$ y la pérdida. (b) Calcula $\partial\ell/\partial\mathbf z$.
(c) Calcula $\mathbf p$ con temperatura $\tau=2$ y $\tau=0.5$. ¿Qué ocurre si $\tau\to0$?
(d) Muestra que restar $\max_kz_k$ a todos los logits no cambia $\mathbf p$.
\end{ejercicio}

\begin{ejercicio}{B18 · XOR: imposibilidad lineal y red de dos capas \dif{2}}
(a) Demuestra que no existen $w_1,w_2,b$ tales que $\mathrm{sign}(w_1x_1+w_2x_2+b)$ reproduzca XOR en $\{0,1\}^2$.
(b) Verifica que $h_1=\mathrm{ReLU}(x_1+x_2)$, $h_2=\mathrm{ReLU}(x_1+x_2-1)$, $\hat y=h_1-2h_2$ calcula XOR.
(c) Cuenta los parámetros de esa red y explica qué hace la capa oculta geométricamente.
\end{ejercicio}

\subsection{Aprendizaje no supervisado y vecinos}

\begin{ejercicio}{B19 · Distancias y $k$-NN \dif{2}\fuente{Clase 29 oct}}
Consulta $q=(3,3)$. Ejemplos etiquetados: $P_1=(6,3)\,[+]$, $P_2=(5,5)\,[-]$, $P_3=(-1,4)\,[+]$, $P_4=(1,-1)\,[-]$.
(a) Calcula $L_1$, $L_2$ y $L_\infty$ de $q$ a cada punto. (b) Clasifica $q$ con 1-NN y 3-NN para cada distancia.
(c) Calcula la distancia de Hamming entre (rojo, grande, metal, sí) y (rojo, chico, plástico, sí).
\end{ejercicio}

\begin{ejercicio}{B20 · $k$-means a mano \dif{2}}
Puntos $P_1(1,1)$, $P_2(2,1)$, $P_3(1,3)$, $P_4(5,4)$, $P_5(6,5)$, $P_6(5,6)$; $K=2$; centroides iniciales $\mu_1=(1,1)$, $\mu_2=(1,3)$.
Itera (asignación con distancia euclidiana al cuadrado, luego actualización) hasta converger. Reporta en cada iteración
la tabla de distancias, la asignación, los nuevos centroides y $J$.
\end{ejercicio}

\begin{ejercicio}{B21 · Silueta \dif{2}}
Con el agrupamiento final de B20, calcula $a(i)$, $b(i)$ y $s(i)$ para $P_3$ y para $P_1$. ¿Cuál está mejor agrupado?
\end{ejercicio}

\begin{ejercicio}{B22 · Agrupamiento jerárquico \dif{2}\fuente{Clase 29 oct, diap.\ 21–27}}
Matriz de distancias:
\begin{center}\small
\begin{tabular}{c|ccccc}
 & A & B & C & D & E\\\hline
A & 0 & 2 & 6 & 10 & 9\\
B & 2 & 0 & 5 & 9 & 8\\
C & 6 & 5 & 0 & 4 & 5\\
D & 10 & 9 & 4 & 0 & 3\\
E & 9 & 8 & 5 & 3 & 0
\end{tabular}
\end{center}
Ejecuta el aglomerativo con \emph{linkage} simple, completo y promedio; da el orden de uniones y sus alturas, y dibuja los
dendrogramas. Si cortas a altura 4.2, ¿cuántos grupos da cada uno?
\end{ejercicio}

\subsection{Aprendizaje por refuerzo}

\begin{ejercicio}{B23 · Retornos, Monte Carlo y TD(0) \dif{2}\fuente{Clase 5 nov}}
$\gamma=0.5$. Episodio 1: $A\xrightarrow{+1}B\xrightarrow{+2}A\xrightarrow{-1}B\xrightarrow{+4}T$. Episodio 2: $B\xrightarrow{+3}A\xrightarrow{+1}T$ ($T$ terminal).
(a) Calcula el retorno desde cada instante (hacia atrás).
(b) Estima $V(A)$ y $V(B)$ con MC de primera visita y de todas las visitas.
(c) Aplica TD(0) con $\alpha=0.5$ y $V_0=0$, en orden, a todas las transiciones de ambos episodios.
\end{ejercicio}

\begin{ejercicio}{B24 · Evaluación exacta de una política \dif{2}}
\begin{center}\figCorredor\end{center}
Corredor con estados $1,2,3$ y terminales $L$, $R$. Acciones izq/der: con probabilidad $0.8$ se mueve en la dirección
pedida y con $0.2$ en la opuesta. Entrar a $L$ da $+2$, entrar a $R$ da $+10$, cualquier otro movimiento $-1$; $\gamma=0.9$.
(a) Escribe las ecuaciones de Bellman para $\pi=$ «siempre derecha» y resuélvelas.
(b) Calcula $Q^\pi(s,\text{izq})$ y $Q^\pi(s,\text{der})$. ¿Mejora la política algún estado? ¿Qué concluyes?
\end{ejercicio}

\begin{ejercicio}{B25 · Iteración de valores en el corredor \dif{2}}
Mismo corredor. Desde $V_0=0$, ejecuta tres iteraciones de valores mostrando los seis valores $Q_k(s,a)$ por iteración.
Extrae la política voraz tras cada iteración. ¿Por qué en las primeras iteraciones el estado 1 prefiere ir a la izquierda?
\end{ejercicio}

\begin{ejercicio}{B26 · Q-learning en el automóvil \dif{2}\fuente{Clase 3 y 5 nov}}
\begin{center}\tabAuto\end{center}
$\alpha=0.5$, $\gamma=0.9$, $Q_0=0$. Se observan, en orden:
$(C,\text{F},+2,W)$, $(W,\text{S},+1,C)$, $(C,\text{F},+2,C)$, $(C,\text{S},+1,C)$, $(C,\text{F},+2,W)$, $(W,\text{F},-10,O)$.
(a) Actualiza la tabla Q tras cada transición (objetivo, valor anterior, valor nuevo). (b) Da la política voraz final.
(c) Compárala con la política óptima del curso. ¿Qué haría falta para que los valores converjan a $Q^*$?
\end{ejercicio}

\begin{ejercicio}{B27 · $\varepsilon$-voraz, SARSA y Q-learning \dif{2}}
(a) Con $Q(s,\cdot)=(3,5,5,1)$ y $\varepsilon=0.2$, da $\mu(a\mid s)$ para las cuatro acciones.
(b) $Q(s,a)=2$, $r=0$, $Q(s',a_1)=6$, $Q(s',a_2)=-2$, $\gamma=0.9$, $\alpha=0.5$. Por exploración, la siguiente acción
realmente tomada es $a_2$. Calcula la actualización de Q-learning y la de SARSA. ¿Cuál aprende el valor de la política
voraz y cuál el de la política que explora?
\end{ejercicio}

\begin{ejercicio}{B28 · Q-learning con aproximación lineal \dif{2}\fuente{Clase 5 nov, diap.\ 15–16}}
$\hat Q(s,a)=\mathbf w^\mathsf T\phi(s,a)$ con $\mathbf w=(1,2,-1)$. Transición: $\phi(s,a)=(1,1,2)$, $r=3$; en $s'$ hay dos
acciones con $\phi=(1,0,1)$ y $(1,2,0)$. $\gamma=0.9$, $\alpha=0.1$.
(a) Calcula $\hat Q(s,a)$, el objetivo y $\delta$. (b) Actualiza $\mathbf w$. (c) ¿Cuánto cambian los valores de las acciones
de $s'$ aunque no se hayan visitado? ¿Por qué se llama semigradiente?
\end{ejercicio}

\subsection{Aprendizaje profundo y evaluación}

\begin{ejercicio}{B29 · Retropropagación en una red 2–2–1 \dif{2}\fuente{Clase 10 nov, estilo}}
$\mathbf x=(1,2)$, $y=1$. Capa oculta ReLU con $W_1=\begin{pmatrix}1&-1\\0.5&0.5\end{pmatrix}$, $\mathbf b_1=(0,-1)$;
salida sigmoide $p=\sigma(\mathbf w_2^\mathsf T\mathbf h+b_2)$, $\mathbf w_2=(2,-1)$, $b_2=0.5$; pérdida entropía cruzada binaria.
(a) Forward: $\mathbf z_1,\mathbf h,z_2,p,\ell$. (b) Backward: $\delta_2$, $\partial\ell/\partial\mathbf w_2$, $\partial\ell/\partial b_2$, $\boldsymbol\delta_1$,
$\partial\ell/\partial W_1$, $\partial\ell/\partial\mathbf b_1$. (c) Actualiza con $\eta=0.5$ y recalcula $p$ y $\ell$.
(d) ¿Por qué la primera fila de $W_1$ no cambia?
\end{ejercicio}

\begin{ejercicio}{B30 · Conteo de parámetros y composición \dif{2}}
(a) Cuenta los parámetros de una red $4\to8\to8\to3$. (b) ¿Cuántos parámetros tiene una tabla de \emph{embeddings} con
$V=32\,000$ y $d=256$? (c) Demuestra que $W_2(W_1\mathbf x+\mathbf b_1)+\mathbf b_2$ es afín. (d) Si una red $2\to3\to2$ usa ReLU,
¿cuántos parámetros tiene y cuántas regiones lineales puede tener como máximo su capa oculta en el plano? (cota: $\le 7$).
\end{ejercicio}

\begin{ejercicio}{B31 · Atención causal a mano \dif{2}\fuente{Clase 10 nov, §12}}
Tres posiciones, $d_k=2$: $\mathbf q_1=(1,0)$, $\mathbf q_2=(0,1)$, $\mathbf q_3=(1,1)$; $\mathbf k_1=(1,0)$, $\mathbf k_2=(1,1)$, $\mathbf k_3=(0,1)$;
$\mathbf v_1=(1,0)$, $\mathbf v_2=(0,1)$, $\mathbf v_3=(1,1)$.
(a) Calcula $S=QK^\mathsf T/\sqrt2$. (b) Aplica la máscara causal y softmax por fila. (c) Calcula $O=AV$.
(d) ¿Qué cambiaría en la fila 2 sin máscara?
\end{ejercicio}

\begin{ejercicio}{B32 · Matriz de confusión y umbral \dif{2}}
Puntajes y etiquetas de 10 casos: $(0.95,1)$, $(0.85,1)$, $(0.80,0)$, $(0.70,1)$, $(0.55,0)$, $(0.45,1)$, $(0.40,0)$, $(0.30,0)$, $(0.20,1)$, $(0.10,0)$.
Para umbrales $0.5$, $0.35$ y $0.75$ calcula $TP,FP,FN,TN$, exactitud, precisión, recall, especificidad y $F_1$.
¿Qué umbral elegirías para detectar una enfermedad grave? ¿Y para un filtro de spam?
\end{ejercicio}

\subsection{IA en el contexto social}

\begin{ejercicio}{B33 · Métricas de equidad por grupo \dif{2}\fuente{Clase 12 nov}}
Un sistema de crédito ($\hat Y=1$: aprobar; $Y=1$: pagaría):
\begin{center}\small
\begin{tabular}{c|cccc}
Grupo & $TP$ & $FN$ & $FP$ & $TN$\\\hline
A & 40 & 10 & 10 & 40\\
B & 15 & 15 & 5 & 65
\end{tabular}
\end{center}
(a) Calcula por grupo: tasa base, tasa de aprobación, TPR, FPR, PPV y exactitud.
(b) ¿Se cumple paridad demográfica? ¿Igualdad de oportunidades? ¿\emph{Equalized odds}?
(c) La exactitud es igual en ambos grupos. ¿Basta eso para decir que el sistema es justo?
\end{ejercicio}

\begin{ejercicio}{B34 · Energía y efecto rebote \dif{2}\fuente{Clase 17 nov}}
(a) ¿Cuántos kWh consume un LED de 10 W en 100 h? ¿Y 100\,000 GPU de 700 W en un día?
(b) Un agente hace 25 llamadas de 0.3 Wh por tarea y acierta en el 75\,\% de las tareas. ¿Cuál es la energía por tarea
completada correctamente? (c) Una mejora reduce la energía por llamada al 40\,\% pero los usos se triplican.
¿Qué ocurre con el consumo total?
\end{ejercicio}

\newpage
% ======================= PARTE C =======================
\section{Parte C — Problemas tipo examen \dif{3}}

\begin{ejercicio}{C1 · Diagnóstico de un servidor: modelar, d-separar y eliminar \dif{3}}
Un servidor web puede sufrir un ataque ($A$) o recibir una actualización defectuosa ($U$); ambos pueden provocar CPU alta ($C$).
La CPU alta dispara una alerta del monitor ($M$) y quejas de usuarios ($Q$).
\begin{center}\figSrv\qquad\tabSrv\end{center}
(a) \pts{4} Escribe la factorización, cuenta parámetros y compara con la conjunta.
(b) \pts{5} Decide, camino por camino: $A\indep U$; $A\indep U\mid M$; $M\indep Q$; $M\indep Q\mid C$; $A\indep Q\mid C$.
(c) \pts{8} Calcula $P(A\mid+m)$ por eliminación de variables (orden $C$, $U$; descarta $Q$). Muestra cada factor.
(d) \pts{5} Calcula $P(A\mid+m,+u)$. Explica qué ocurrió.
(e) \pts{3} Calcula $P(A\mid+m,+q)$. ¿Por qué la segunda señal aumenta la creencia de ataque menos de lo que podría esperarse?
\end{ejercicio}

\begin{ejercicio}{C2 · Enumeración frente a eliminación en la red del estudiante \dif{3}}
Usa la red de B4. Consulta: $P(I\mid+r,+e)$.
(a) \pts{8} Resuelve por enumeración: lista los 8 mundos compatibles con su probabilidad y normaliza.
(b) \pts{8} Resuelve por eliminación con orden $D,C$, reutilizando los factores de B8.
(c) \pts{4} ¿Cuántos productos de 5 factores hace la enumeración ingenua? ¿Cuántas entradas tiene el factor más grande de la eliminación?
(d) \pts{5} Calcula $P(+i\mid+c)$, $P(+i\mid+c,+d)$ y $P(+i\mid+c,-d)$. Interpreta: ¿por qué saber que el curso era difícil
\emph{aumenta} la creencia en la inteligencia? ¿Contradice esto el «explicar a la otra causa»?
\end{ejercicio}

\begin{ejercicio}{C3 · Comparar métodos de muestreo \dif{3}}
Red del servidor (C1); consulta $P(+a\mid+m)$ (exacto en C1). La variable $Q$ es irrelevante y se omite.
(a) \pts{6} \emph{Rechazo.} Se generaron 12 muestras previas $(A,U,C,M)$:
\begin{center}\small
\begin{tabular}{llll}
$-a,-u,-c,-m$ & $-a,+u,+c,+m$ & $-a,-u,-c,-m$ & $+a,-u,+c,+m$\\
$-a,-u,+c,-m$ & $-a,-u,-c,+m$ & $-a,-u,-c,-m$ & $-a,+u,-c,-m$\\
$+a,-u,+c,-m$ & $-a,-u,-c,-m$ & $-a,-u,-c,-m$ & $-a,+u,+c,+m$
\end{tabular}
\end{center}
¿Cuántas se aceptan y qué estimación dan? ¿Cuántas aceptadas esperarías de 12 en promedio?
(b) \pts{7} \emph{Verosimilitud con pesos.} Con $M=+m$ fijo se muestrearon $(A,U,C)$:
$(-a,-u,-c)$, $(-a,+u,+c)$, $(+a,-u,+c)$, $(-a,-u,-c)$, $(-a,-u,+c)$, $(-a,-u,-c)$, $(-a,-u,-c)$, $(-a,+u,-c)$.
Calcula los pesos y la estimación.
(c) \pts{6} \emph{Gibbs.} Estado $(-a,+u,+c)$ con $+m$. Calcula $P(A\mid\text{manto})$ y, si después se actualiza $C$ con $A=-a$, $P(C\mid\text{manto})$.
(d) \pts{6} Compara los tres métodos en esta consulta: ¿cuál desperdicia más trabajo y por qué? ¿Cuál preferirías si la evidencia fuese $+m,+q$ y muy rara?
\end{ejercicio}

\begin{ejercicio}{C4 · Clasificador bayesiano ingenuo como red bayesiana \dif{3}\fuente{AIMA 20.2, estilo}}
¿Se juega fútbol ($J$)? Datos:
\begin{center}\small
\begin{tabular}{c|cccccccccc}
\# & 1&2&3&4&5&6&7&8&9&10\\\hline
Cielo & sol&sol&nub&lluvia&lluvia&nub&sol&lluvia&nub&sol\\
Viento & sí&no&no&sí&no&sí&no&sí&no&sí\\
$J$ & no&sí&sí&no&sí&sí&sí&no&sí&no
\end{tabular}
\end{center}
(a) \pts{4} Dibuja la red ($J\to\text{Cielo}$, $J\to\text{Viento}$) y escribe la regla de clasificación.
(b) \pts{7} Estima por MLE $P(J)$, $P(\text{Cielo}\mid J)$ y $P(\text{Viento}\mid J)$.
(c) \pts{7} Clasifica (lluvia, sí), (nub, sí) y (sol, no) con MLE. ¿Qué problema aparece?
(d) \pts{7} Repite (c) con Laplace $\lambda=1$ en las CPT de los atributos ($P(J)$ por MLE). ¿Cambia alguna decisión o solo la confianza?
\end{ejercicio}

\begin{ejercicio}{C5 · EM con dos monedas \dif{3}\fuente{Do \& Batzoglou (2008), clásico}}
Se eligió al azar (prob.\ $1/2$) una de dos monedas $A$, $B$ y se lanzó 10 veces; se repitió 5 veces. Solo se registró el número
de caras: $5,9,8,4,7$. Parámetros iniciales $\theta_A=0.6$, $\theta_B=0.5$.
(a) \pts{4} Plantea la red bayesiana (moneda oculta $\to$ resultados) y la verosimilitud marginal.
(b) \pts{10} Paso E: calcula $w_j=P(A\mid h_j)$ para cada serie (usa $\theta^h(1-\theta)^{10-h}$).
(c) \pts{6} Paso M: calcula las caras y cruces esperadas de cada moneda y reestima $\theta_A,\theta_B$.
(d) \pts{5} Si se conociera qué moneda se usó en cada serie, ¿cómo se estimaría? ¿Qué pasa si se intercambian las etiquetas $A$ y $B$?
\end{ejercicio}

\begin{ejercicio}{C6 · Regresión polinomial y generalización \dif{3}}
Entrenamiento: $x=-2,-1,0,1,2$ con $y=4.2,0.8,0.1,1.1,3.8$. Prueba: $(3,9.1)$ y $(-3,8.8)$.
(a) \pts{6} Ajusta $\hat y=\theta_0+\theta_1x$ por ecuaciones normales. Calcula ECM de entrenamiento y de prueba.
(b) \pts{9} Ajusta $\hat y=\theta_0+\theta_1x+\theta_2x^2$. (Pista: $X^\mathsf TX$ tiene muchos ceros por la simetría de los $x$.) Calcula ambos ECM.
(c) \pts{5} Un polinomio de grado 4 interpola los cinco puntos con $\theta=(0.1,\,0.2333,\,0.8083,\,-0.0833,\,0.0417)$.
Verifica su ECM de entrenamiento y calcula sus predicciones de prueba.
(d) \pts{5} Ordena los tres modelos por error de entrenamiento y por error de prueba; nombra subajuste y sobreajuste.
\end{ejercicio}

\begin{ejercicio}{C7 · $k$-means y la escala de las características \dif{3}}
Clientes (ingreso en miles, visitas por mes): $C_1(10,1)$, $C_2(12,9)$, $C_3(30,2)$, $C_4(32,8)$. $K=2$, centroides iniciales $C_1$ y $C_2$.
(a) \pts{8} Ejecuta $k$-means sobre los datos crudos hasta converger; reporta $J$.
(b) \pts{8} Divide el ingreso entre 10 y repite. ¿Qué grupos obtienes ahora?
(c) \pts{5} Calcula la silueta de $C_1$ en cada solución (en su propia escala).
(d) \pts{4} ¿Cuál agrupamiento es «correcto»? ¿Qué decide la respuesta?
\end{ejercicio}

\begin{ejercicio}{C8 · Robot de reciclaje: modelar y resolver el MDP \dif{3}\fuente{Sutton \& Barto, ej.\ 3.3, adaptado}}
Un robot recolecta latas. Su batería está alta ($H$) o baja ($L$). En $H$ puede \emph{buscar} o \emph{esperar}; en $L$ puede \emph{buscar},
\emph{esperar} o \emph{recargar}.
\begin{itemize}
\item Buscar en $H$: recompensa $+2$; queda en $H$ con prob.\ $0.7$ y pasa a $L$ con $0.3$.
\item Buscar en $L$: con prob.\ $0.6$ recompensa $+2$ y sigue en $L$; con $0.4$ se agota, lo rescatan (recompensa $-3$) y vuelve a $H$.
\item Esperar: recompensa $+1$ y no cambia de estado. Recargar (solo en $L$): recompensa $0$ y pasa a $H$. $\gamma=0.9$.
\end{itemize}
(a) \pts{5} Escribe el MDP: $S$, $A(s)$, $T$ y $R$ en una tabla; dibuja el diagrama de transiciones.
(b) \pts{10} Iteración de políticas desde $\pi_0=(\text{esperar},\text{esperar})$: evalúa exactamente, mejora y repite hasta que no cambie.
(c) \pts{6} Cuatro iteraciones de valores desde $V_0=0$. ¿Qué política voraz da $V_4$?
(d) \pts{4} ¿Cómo cambiaría la política óptima si la penalización de rescate fuera $-20$? (Razona sin resolver todo.)
\end{ejercicio}

\begin{ejercicio}{C9 · Q-learning en el corredor y comparación con el óptimo \dif{3}}
Corredor de B24 ($\gamma=0.9$), Q-learning con $\alpha=0.5$, $Q_0=0$. Episodios observados (estado, acción, recompensa, siguiente):
\begin{enumerate}[label=Ep.\,\arabic*:,leftmargin=3em]
\item $(2,\text{der},-1,3)$, $(3,\text{der},+10,R)$
\item $(1,\text{der},-1,2)$, $(2,\text{der},-1,3)$, $(3,\text{izq},-1,2)$, $(2,\text{der},-1,3)$, $(3,\text{der},+10,R)$
\item $(2,\text{izq},-1,1)$, $(1,\text{izq},+2,L)$
\item $(1,\text{der},-1,2)$, $(2,\text{der},-1,3)$, $(3,\text{der},+10,R)$
\end{enumerate}
(a) \pts{12} Actualiza la tabla Q tras cada transición. (b) \pts{4} Da la política voraz final.
(c) \pts{5} Compara con $Q^*$ (usa B24): ¿qué estado tiene la acción equivocada y por qué?
(d) \pts{4} En el episodio 2 el agente eligió «izq» en el estado 3 (exploración). Explica por qué Q-learning sigue
aprendiendo el valor de la política voraz y qué cambiaría con SARSA. Si hubiera pedido «der» y el ruido lo hubiera llevado a 2,
¿qué tupla se registraría y qué entrada se actualizaría?
\end{ejercicio}

\begin{ejercicio}{C10 · Monte Carlo frente a TD: el ejemplo A–B \dif{3}\fuente{Sutton \& Barto, ej.\ 6.4}}
$\gamma=1$. Ocho episodios: $A\xrightarrow{0}B\xrightarrow{0}T$; seis veces $B\xrightarrow{1}T$; una vez $B\xrightarrow{0}T$.
(a) \pts{6} Estima $V(A)$ y $V(B)$ con Monte Carlo.
(b) \pts{8} Estima $V(A)$ y $V(B)$ con TD en modo \emph{batch} (repitiendo los episodios hasta converger). Pista: busca el punto fijo.
(c) \pts{6} ¿Cuál estimación de $V(A)$ prefieres y por qué? Relaciona con «modelo de máxima verosimilitud» y con el error en los datos.
(d) \pts{5} Aplica una sola pasada de TD(0) en línea con $\alpha=0.5$, $V_0=0$, en el orden dado.
\end{ejercicio}

\begin{ejercicio}{C11 · Retropropagación con softmax \dif{3}}
$\mathbf x=(1,-1)$, clase correcta 2 ($\mathbf t=(0,1)$). Oculta ReLU: $W_1=\begin{pmatrix}2&1\\1&-1\end{pmatrix}$, $\mathbf b_1=(0,-3)$.
Salida: $\mathbf z_2=W_2\mathbf h+\mathbf b_2$, $W_2=\begin{pmatrix}1&-1\\0.5&1\end{pmatrix}$, $\mathbf b_2=(0,0)$, softmax y entropía cruzada.
(a) \pts{6} Forward completo y pérdida. (b) \pts{10} Backward: $\boldsymbol\delta_2$, $\partial\ell/\partial W_2$, $W_2^\mathsf T\boldsymbol\delta_2$, $\boldsymbol\delta_1$, $\partial\ell/\partial W_1$.
(c) \pts{5} Actualiza con $\eta=1$ y recalcula la pérdida. (d) \pts{4} ¿Qué parámetros no se actualizan y por qué? ¿Qué es una neurona muerta?
\end{ejercicio}

\begin{ejercicio}{C12 · Un transformer diminuto \dif{3}}
Vocabulario $\{\text{el},\text{gato},\text{come},\text{pez},\langle/s\rangle\}$ (en ese orden). Se entrena con la secuencia
«$\langle s\rangle$ el gato come $\langle/s\rangle$». Los logits por posición son:
\begin{center}\small
\begin{tabular}{l|ccccc}
Prefijo & el & gato & come & pez & $\langle/s\rangle$\\\hline
$\langle s\rangle$ & 2 & 0 & 0 & 0 & $-1$\\
$\langle s\rangle$ el & 0 & 2 & 0 & 1 & 0\\
$\langle s\rangle$ el gato & 1 & 0 & 2 & 0 & 0\\
$\langle s\rangle$ el gato come & 0 & 0 & 1 & 0 & 3
\end{tabular}
\end{center}
(a) \pts{8} Calcula la probabilidad del token correcto en cada posición, la pérdida media y la perplejidad.
(b) \pts{4} Para el tercer prefijo, calcula la distribución con $\tau=0.5$ y $\tau=2$.
(c) \pts{9} Atención causal de una cabeza con $Q=\begin{pmatrix}1&1\\0&2\\2&0\end{pmatrix}$, $K=\begin{pmatrix}1&0\\0&1\\1&1\end{pmatrix}$,
$V=\begin{pmatrix}2&0\\0&2\\1&1\end{pmatrix}$: calcula $S$, $A$ y $O$.
(d) \pts{4} Con $V=50\,000$, $d=512$ y una cabeza de $d_k=d_v=64$, cuenta los parámetros de la tabla de embeddings y de $W_Q,W_K,W_V,W_O$.
\end{ejercicio}

\begin{ejercicio}{C13 · Auditoría de un sistema de crédito \dif{3}}
Resultados en dos grupos de 100 personas:
\begin{center}\small
\begin{tabular}{c|cccc}
Grupo & $TP$ & $FN$ & $FP$ & $TN$\\\hline
A & 48 & 12 & 8 & 32\\
B & 18 & 12 & 2 & 68
\end{tabular}
\end{center}
(a) \pts{8} Calcula exactitud global y, por grupo, tasa base, aprobación, TPR, FPR, PPV y exactitud.
(b) \pts{6} ¿Qué criterios de equidad se violan? ¿Quién paga el costo de los errores en cada grupo?
(c) \pts{4} ¿Por qué no se pueden igualar a la vez TPR, FPR y PPV aquí? ¿Qué es el problema de las etiquetas selectivas?
(d) \pts{4} ¿En qué nivel de riesgo del Reglamento de IA de la UE cae este sistema y qué obligaciones implica?
(e) \pts{3} El sistema usa un agente que hace 40 llamadas de 0.3 Wh por solicitud y resuelve bien el 80\,\%. Calcula la energía
por solicitud bien resuelta y para 10\,000 solicitudes diarias.
\end{ejercicio}

\newpage
% ======================= PARTE D =======================
\section{Parte D — Exámenes simulados (2 horas cada uno)}

\begin{nota}
Resuelve cada simulacro de corrido, con cronómetro y con tu hoja A4. Cuatro preguntas de 25 puntos. Califícate con las
respuestas: da la mitad de los puntos de un inciso si el procedimiento es correcto y solo falla la aritmética.
\end{nota}

\subsection*{Simulacro D1}
\addcontentsline{toc}{subsection}{Simulacro D1}

\begin{ejercicio}{D1.1 · El auto no arranca (25 pts) \dif{3}}
\begin{center}\figCar\qquad
{\small\begin{tabular}{l}
$B$: batería bien ($P(+b)=0.9$); $F$: hay gasolina ($P(+f)=0.95$)\\
$S$: arranca: $P(+s\mid+b,+f)=0.99$, $P(+s\mid+b,-f)=0$,\\
\hspace{3.5em}$P(+s\mid-b,+f)=0.05$, $P(+s\mid-b,-f)=0$\\
$L$: luces encienden: $P(+l\mid+b)=0.95$, $P(+l\mid-b)=0.1$
\end{tabular}}\end{center}
(a) \pts{4} Factorización y número de parámetros. (b) \pts{6} ¿$B\indep F$? ¿$B\indep F\mid S$? ¿$L\indep F$? ¿$L\indep F\mid S$? Justifica.
(c) \pts{8} Calcula $P(-b\mid-s)$ y $P(-f\mid-s)$ por enumeración.
(d) \pts{7} Calcula $P(-b\mid-s,+l)$ y $P(-f\mid-s,+l)$. Explica en una línea cómo las luces cambian el diagnóstico.
\end{ejercicio}

\begin{ejercicio}{D1.2 · Aprender y muestrear la red del auto (25 pts) \dif{3}}
Registros $(B,F,S,L)$ con $1=$ verdadero: $(1,1,1,1)$, $(1,1,1,1)$, $(1,0,0,1)$, $(0,1,0,0)$, $(1,1,1,1)$, $(0,1,1,0)$, $(1,1,1,0)$, $(0,1,0,0)$.
(a) \pts{10} Estima por MLE y con Laplace ($\lambda=1$) las cuatro filas de $P(+s\mid B,F)$ y las dos de $P(+l\mid B)$. ¿Qué fila es indefinida con MLE?
(b) \pts{10} Con las CPT \emph{originales} de D1.1, verosimilitud con pesos para $P(+b\mid-s)$: se muestrearon $(B,F)=(1,1),(0,1),(1,0),(1,1),(0,1),(1,1)$.
Calcula pesos y estimación; compara con el valor exacto.
(c) \pts{5} ¿Por qué rechazo sería mala idea para $P(B\mid-s)$? Estima cuántas muestras previas necesitarías para aceptar 100.
\end{ejercicio}

\begin{ejercicio}{D1.3 · Inventario como MDP (25 pts) \dif{3}}
Una tienda tiene $s\in\{0,1,2\}$ unidades. Cada día decide pedir $a\in\{0,1\}$ unidades (sin superar 2); el pedido llega antes de la
demanda. La demanda es 1 o 0 con probabilidad $0.5$ cada una. Vender una unidad da $+4$, pedir cuesta $1$ por unidad y cada unidad
que queda al final del día cuesta $0.5$. $\gamma=0.8$.
(a) \pts{6} Escribe $T(s,a,s')$ y $R(s,a,s')$ para los 5 pares $(s,a)$ válidos.
(b) \pts{9} Tres iteraciones de valores desde $V_0=0$ y la política voraz de $V_3$.
(c) \pts{6} Evalúa exactamente $\pi=$ «pedir si $s<2$» y $\pi'=$ «pedir solo si $s=0$». ¿Cuál es mejor? ¿Es $\pi'$ óptima? (verifica con un paso de mejora).
(d) \pts{4} ¿Qué cambia si $\gamma=0$?
\end{ejercicio}

\begin{ejercicio}{D1.4 · Q-learning tabular y aproximado (25 pts) \dif{3}}
(a) \pts{12} Tabla inicial: $Q(s_1,a)=1$, $Q(s_1,b)=2$, $Q(s_2,a)=4$, $Q(s_2,b)=0$, $Q(s_3,a)=-1$, $Q(s_3,b)=3$; $\alpha=0.25$, $\gamma=0.8$.
Transiciones: $(s_1,b,+1,s_2)$, $(s_2,a,-2,s_3)$, $(s_3,b,+5,T)$, $(s_1,a,0,s_3)$. Actualiza en orden.
(b) \pts{5} Da la política $\varepsilon$-voraz ($\varepsilon=0.1$) en $s_1$ antes y después.
(c) \pts{8} Aproximación lineal: $\mathbf w=(0.5,1,-2)$, $\phi(s,a)=(1,3,1)$, $r=-1$, acciones en $s'$ con $\phi=(1,1,0)$ y $(1,2,1)$,
$\gamma=0.8$, $\alpha=0.1$. Calcula $\delta$ y el nuevo $\mathbf w$.
\end{ejercicio}

\subsection*{Simulacro D2}
\addcontentsline{toc}{subsection}{Simulacro D2}

\begin{ejercicio}{D2.1 · Lluvia, aspersor y tráfico (25 pts) \dif{3}}
\begin{center}\figDdos\qquad
{\small\begin{tabular}{l}
$P(+n)=0.4$;\ $P(+r\mid+n)=0.7$, $P(+r\mid-n)=0.1$\\
$P(+s\mid+n)=0.2$, $P(+s\mid-n)=0.6$\\
$P(+w\mid R,S)$: $(+,+)\,0.95$, $(+,-)\,0.85$, $(-,+)\,0.8$, $(-,-)\,0.05$\\
$P(+t\mid+r)=0.6$, $P(+t\mid-r)=0.2$
\end{tabular}}\end{center}
(a) \pts{8} d-separación: $S\indep T$; $S\indep T\mid R$; $S\indep T\mid W$; $S\indep T\mid R,W$; $N\indep W\mid R,S$; $T\indep W\mid R$.
(b) \pts{10} $P(R\mid+w)$ por eliminación con orden $N,S$ (descarta $T$ justificando). Muestra cada factor.
(c) \pts{7} $P(R\mid+w,+s)$. Compara con (b) y explica.
\end{ejercicio}

\begin{ejercicio}{D2.2 · Regresión, logística y métricas (25 pts) \dif{3}}
(a) \pts{9} Ajusta $\hat y=\theta_0+\theta_1x$ a $(1,2),(2,3),(4,6),(5,9)$; da SSE y la predicción en $x=6$.
(b) \pts{10} Logística con $\theta=(\theta_0,\theta_1)=(0,0)$ y datos $x=1,3$ (clase 1), $x=-1,-2$ (clase 0). Calcula la pérdida, el gradiente,
un paso con $\alpha=0.5$ y la nueva pérdida. ¿Qué pasa con $\theta$ si se sigue iterando? ¿Por qué?
(c) \pts{6} Con $TP=30$, $FP=10$, $FN=20$, $TN=140$ calcula exactitud, precisión, recall, especificidad y $F_1$.
\end{ejercicio}

\begin{ejercicio}{D2.3 · Agrupamiento (25 pts) \dif{3}}
Puntos $A(0,0)$, $B(0,2)$, $C(3,0)$, $D(6,6)$, $E(7,6)$, $F(9,9)$.
(a) \pts{12} $k$-means ($K=2$, euclidiana al cuadrado) con centroides iniciales $A$ y $B$ hasta converger; reporta $J$. Repite con $A$ y $D$. Comenta.
(b) \pts{10} Con distancia Manhattan, aglomerativo con \emph{linkage} simple y completo: orden de uniones y alturas.
(c) \pts{3} ¿Cuántos grupos sugiere cada dendrograma? ¿Por qué?
\end{ejercicio}

\begin{ejercicio}{D2.4 · Red neuronal, atención y equidad (25 pts) \dif{3}}
(a) \pts{12} $\mathbf x=(2,1)$, $y=0$. $z_1=\mathbf w_1^\mathsf T\mathbf x+b_1$, $h=\mathrm{ReLU}(z_1)$, $p=\sigma(w_2h+b_2)$, entropía cruzada;
$\mathbf w_1=(1,-1)$, $b_1=0$, $w_2=2$, $b_2=-1$. Forward, backward, actualización con $\eta=0.5$ y nueva pérdida.
(b) \pts{7} Consulta $\mathbf q=(1,2)$, claves $(2,0),(0,1),(1,1)$, valores $(1,0),(0,1),(2,2)$, $d_k=2$. Calcula pesos y salida sin máscara.
Si la consulta está en la posición 2 con máscara causal, ¿qué cambia?
(c) \pts{6} Grupos H: $TP=30,FN=10,FP=15,TN=45$; M: $TP=20,FN=20,FP=5,TN=55$. Compara aprobación, TPR, FPR, PPV y exactitud. ¿Qué criterio se viola?
\end{ejercicio}

\end{document}
